\chapter{Contexte et présentation de la mission}
\chaptermark{Contexte et présentation}
\label{chap:contexte}

\section{Cadre du projet}

\subsection{Organisme d'accueil}

Le stage s'est déroulé à l'Université de Mons (UMONS), au sein du service
Informatique, Logiciel et Intelligence Artificielle (ILIA) de la Faculté
Polytechnique, implanté rue de Houdain à Mons. Ce service mène des activités
d'enseignement et de recherche en informatique. Ses travaux couvrent notamment
l'intelligence artificielle, la vision par ordinateur, la gestion de données
issues de l'Internet des objets ainsi que les environnements Cloud et
Edge~Computing. Ses enseignements incluent également le génie logiciel, la
programmation embarquée, l'IoT et la cybersécurité.

L'encadrement du stage a été assuré par M.~Saïd Mahmoudi, membre du service
ILIA. Ce cadre associait ainsi les compétences du service
en intelligence artificielle, systèmes distribués et IoT aux problématiques de
sécurité traitées par LANCE.

\subsection{Organisation technique}

Le développement s'appuie sur deux espaces aux responsabilités distinctes. Le
dépôt canonique rassemble le code versionné de l'agent, de l'\gls{api}, de
l'interface, du benchmark et de l'infrastructure. Il constitue la référence pour
les fonctionnalités, les configurations, les tests et la documentation associés
à une version donnée. L'espace d'entraînement accueille les traces préparées,
les jeux de données, les checkpoints et les adaptateurs produits pendant les
expériences.

Cette séparation évite de placer dans le dépôt des artefacts volumineux ou des
données sensibles, tout en maintenant la traçabilité expérimentale. Une
campagne doit pouvoir être reliée à une configuration du modèle, à une version
du pipeline, à un jeu de données identifié et à un protocole d'évaluation.
Les scripts et configurations restent donc versionnés, tandis que données et
poids sont conservés dans l'espace prévu à cet effet.

\section{Système existant}

LANCE associe trois ensembles fonctionnels~:

\begin{description}
  \item[Agent d'audit] L'agent conduit une campagne en six phases, de l'analyse
  initiale du réseau à la synthèse. Chaque phase reçoit un contexte défini,
  dispose d'outils autorisés et produit un artefact structuré pour la suivante.

  \item[Pipeline d'exécution] Le pipeline organise les transitions, parallélise
  certaines analyses, valide les sorties et applique les règles de sécurité.
  Il transforme les décisions probabilistes en actions observables et
  contrôlables.

  \item[Harness expérimental] Le benchmark décrit des réseaux
  \gls{iot}/\gls{ot} reproductibles à partir de topologies, de vulnérabilités
  injectées et de postures de sécurité. Il déploie les scénarios, conserve une
  vérité terrain et calcule les mesures d'évaluation.
\end{description}

Une \gls{api} assure le lancement et le suivi des campagnes, tandis qu'une
interface présente la topologie, les événements et les artefacts. La chaîne
d'apprentissage mobilise PyTorch, Transformers, PEFT, bitsandbytes et TRL ; le
déploiement des environnements repose sur des descriptions déclaratives et une
automatisation de l'infrastructure. Les résultats sont enregistrés dans des
formats structurés afin d'être analysés indépendamment de la génération du
modèle.

\section{Limites initiales}

\begin{description}
  \item[Comportement trop généraliste] Un même modèle assumait des tâches aux
  contraintes différentes. Les rôles dépendaient principalement des prompts,
  sans garantir la conformité des appels d'outils ni la stabilité du
  comportement.

  \item[Spécialisation incomplète] Il manquait un processus continu pour
  nettoyer les traces, produire des exemples, adapter le modèle avec des
  ressources limitées, servir plusieurs spécialités et les sélectionner pendant
  une campagne.

  \item[Pipeline insuffisamment contractuel] Une sortie syntaxiquement valide
  pouvait encore porter une hypothèse non vérifiée ou être transmise à la phase
  suivante. Les interruptions et erreurs d'outils devaient aussi être
  représentées sans ambiguïté.

  \item[Mesure partielle] La comparaison à la vérité terrain ne caractérisait
  pas assez finement le comportement de l'agent. Une évaluation sémantique seule
  aurait, à l'inverse, rendu les scores sensibles au modèle juge et à la
  formulation.

  \item[Contexte volumineux] Les topologies, sorties de commandes et observations
  répétées pouvaient saturer la fenêtre du modèle, alourdir le traitement et masquer
  des informations nécessaires à la décision.

  \item[Priorisation topologique limitée] La représentation du réseau décrivait
  les hôtes et leurs relations, mais évaluait peu l'effet d'une nouvelle
  découverte sur les chemins menant à une cible.
\end{description}

\section{Mission principale~: spécialisation \gls{hmoe} agentique}

La mission prioritaire était de construire une architecture \gls{hmoe} adaptée à
LANCE. Contrairement à un mélange d'experts neuronal qui route les jetons, elle
partage un modèle quantifié entre des adaptateurs LoRA spécialisés par rôle ; le
routeur active l'un d'eux selon la phase ou le rôle demandé.

Le travail couvre la préparation, la validation et la diversification des traces,
l'entraînement \gls{qlora}, le service derrière une \gls{api} commune et le
routage. Il comprend aussi l'intégration des experts à l'agent qui planifie,
appelle les outils, interprète leurs réponses et transmet un état structuré, avec
un routage stable, des changements d'adaptateur sérialisés et une compatibilité
avec l'existant.

\section{Pipeline et harness comme protocole de validation}

La spécialisation ne se juge pas sur la seule perte d'entraînement. Le pipeline
sert de protocole opérationnel : il délimite les rôles, injecte les outils,
valide les schémas et les preuves, maintient la cohérence entre phases et signale
toute campagne incomplète.

Le harness fournit le protocole expérimental complémentaire : il crée des
situations contrôlées, associe les faiblesses attendues à des indicateurs et
compare les artefacts à la vérité terrain. Les scénarios, leur déploiement, le
suivi des campagnes et les métriques permettent ainsi de déterminer si la
spécialisation améliore l'action de l'agent au-delà du style de ses réponses.

\section{Objectifs opérationnels}

La mission se décline en trois objectifs~:

\begin{enumerate}
  \item concevoir la préparation des données, l'entraînement \gls{qlora}, le
  service multi-adaptateurs et le routage des experts ;
  \item consolider le pipeline afin de rendre transitions, appels d'outils et
  preuves vérifiables ;
  \item étendre le harness pour déployer des scénarios reproductibles et mesurer
  le comportement des différentes configurations.
\end{enumerate}

Ces objectifs forment un même système : les experts décident, le pipeline encadre
leur exécution et le harness les compare. Les optimisations de contexte et de
graphe soutiennent cet ensemble sans remplacer la spécialisation agentique.
